% !TeX root = ../../Book2.tex
\section{必要条件}
\label{b2:sec:2302}

限制估计不是单凭 Fourier 变换的有界性就能得到的。曲面上的质量可以集中在小帽，也可以遍布整个球面；前一种输入产生很长的相干管，后一种输入留下缓慢衰减的尾部。本节分别计算这些障碍。计算所得是不等式成立的必要条件，后面的充分估计还需要独立证明。

\subsection{缩放}
\label{b2:subsec:230201}

先直接研究曲面上的积分
\[
 Eg(x)\coloneqq\int_S e^{2\pi\mathrm i x\cdot\xi}
                     g(\xi)\,d\sigma(\xi),
 \quad \|Eg\|_{L^q(\mathbb R^d)}
                 \leq C\|g\|_{L^s(S)}.
\]
\symidx{\(Eg\)：曲面延拓算子}
其中 \(S\) 暂取单位球面，\(1\leq s,q\leq\infty\)。球面测度有限，故
\(g\in L^s(S)\) 总属于 \(L^1(S)\)，积分处处绝对收敛。
这里先把 \(E\) 当作一个明确的积分算子；它与限制估计的对偶等价将在
\cref{b2:prop:23-duality}中证明。对应的限制指数是
\(p=q'\)、\(r=s'\)。

\ProseParagraphBreak
设半径 \(R\) 的球面取原生面积测度，\(g_R\) 是其上的函数，并令
\(\widetilde g(\omega)=g_R(R\omega)\)。直接换元得
\[
 E_Rg_R(x)=R^{d-1}E\widetilde g(Rx),\quad
 \|g_R\|_{L^s(RS)}
       =R^{(d-1)/s}\|\widetilde g\|_{L^s(S)}.
\]
若单位球面的最佳常数为 \(C\)，半径 \(R\) 的最佳常数便为
\[
 C_R=C R^{\,d-1-d/q-(d-1)/s}.
\]
因为输入缩放可逆，这里也是准确的最佳常数关系。
它与\cref{b2:prop:23-surface-scaling}在共轭指数下完全一致。
固定单位球面并不要求这个幂等于零；只有要求一族半径上的常数统一时，才能由 \(R\to0\) 和 \(R\to\infty\) 进一步限制该幂。
本节真正限制固定球面指数的，是球面内部任意小帽的尺度。

\ProseParagraphBreak
还须区分两个空间的测度。曲面有限测度给出
\[
 \|g\|_{s_1}\leq \sigma(S)^{1/s_1-1/s_2}\|g\|_{s_2},
 \quad 1\leq s_1\leq s_2\leq\infty.
\]
环境空间的测度却是无穷大，不能据 \(q\) 的大小直接写出一般
\(L^q(\mathbb R^d)\) 之间的包含。不过延拓函数另有
\(\|Eg\|_\infty\leq\|g\|_1\)。如果已经证明某个输出指数 \(q_0<\infty\)
的界，那么
\[
 \|Eg\|_q\leq
 \|Eg\|_{q_0}^{q_0/q}
       \cdot\|Eg\|_\infty^{\,1-q_0/q},
 \quad q_0\leq q<\infty,
\]
才给出更大输出指数的界。这一步使用的是输出的额外有界性。

\subsection{Knapp 例子}
\label{b2:subsec:230202}

球面小帽在切向上的半径为 \(\delta\)，离开切平面的高度只有
\(O(\delta^2)\)。为了使不同频率的相位保持接近，环境变量的切向分量可以达到
\(\delta^{-1}\)，法向分量则可以达到 \(\delta^{-2}\)。
这就是下面的\term[knapp-lizi]{Knapp 例子}[Knapp example]所用的几何。

\Needspace{6\baselineskip}
\begin{proposition}[帽与管的必要条件]\label{b2:prop:23-knapp}
设 \(d\geq2\)，单位球面 \(S^{d-1}\) 取面积测度，
\(1\leq s,q\leq\infty\)。若其延拓积分对全部
\(g\in L^s(S^{d-1})\) 满足
\(\|Eg\|_q\leq C\|g\|_s\)，则必须有
\[
 \frac{d+1}{q}\leq\frac{d-1}{s'},
 \quad \frac1s+\frac1{s'}=1.
\]
\end{proposition}
\begin{proof}
在北极附近写
\(\xi=(\eta,h(\eta))\)，其中
\(h(\eta)=\sqrt{1-|\eta|^2}\)。取 \(0<\delta<1/4\)，令
\[
 \Theta_\delta=\{(\eta,h(\eta)):|\eta|\leq\delta\},
 \quad g_\delta=\mathbf1_{\Theta_\delta}.
\]
面积因子 \((1-|\eta|^2)^{-1/2}\) 在这些帽上有统一正上下界，所以
\[
 \sigma(\Theta_\delta)\asymp\delta^{d-1},\quad
 \|g_\delta\|_s\asymp\delta^{(d-1)/s}.
\]
当 \(s=\infty\) 时后一范数等于一，仍与所写幂一致。

由 \(h(\eta)-1=-|\eta|^2/2+O(|\eta|^4)\)，存在与 \(\delta\)
无关的 \(C_0\)，使 \(|h(\eta)-1|\leq C_0\delta^2\)。
取一个固定的小 \(c>0\)，在集合
\[
 \mathcal T_\delta=
 \{x:|x'|\leq c\delta^{-1},\ |x_d|\leq c\delta^{-2}\}
\]
上，有
\[
 2\pi|x'\cdot\eta+x_d(h(\eta)-1)|
       \leq 2\pi c(1+C_0).
\]
选 \(c\) 使右侧不超过 \(\pi/3\)。从积分中提出公共因子
\(e^{2\pi\mathrm i x_d}\)，余下被积函数的实部至少为 \(1/2\)，故
\[
 |Eg_\delta(x)|\geq\tfrac12\sigma(\Theta_\delta)
       \gtrsim\delta^{d-1}\quad(x\in\mathcal T_\delta).
\]
管的体积为常数乘 \(\delta^{-(d+1)}\)。因此，包括 \(q=\infty\)
的通常解释在内，
\[
 \|Eg_\delta\|_q\gtrsim
       \delta^{d-1-(d+1)/q}.
\]
假定的估计会要求
\(\delta^{d-1-(d+1)/q}\lesssim\delta^{(d-1)/s}\)。
令 \(\delta\downarrow0\)，左侧幂不能小于右侧幂，整理即得结论。
\end{proof}

同样的构造可以用非负光滑帽代替特征函数：令帽内部保留一个固定比例的核心，振幅在核心上有正下界即可。因此，必要条件并非由输入的不光滑造成。对任意光滑曲面的局部片，减去切平面后只需使用高度的 \(O(\delta^2)\) 上界，也可得到同一个必要条件；若曲率退化，还可能有更强的条件。

\ProseParagraphBreak
当 \(s=2\) 时，上式成为
\[
 q\geq q_*\coloneqq\frac{2(d+1)}{d-1}.
\]
本章将证明等号成立，因此平方曲面输入的最小输出指数已经由一个帽确定。
例如 \(d=3\) 时帽面积约为 \(\delta^2\)，管体积约为 \(\delta^{-4}\)，
比较 \(\delta^{2-4/q}\) 与 \(\delta\) 就得到 \(q\geq4\)。
若 \(s=1\)，必要条件要求 \(q=\infty\)，正好对应绝对积分界。

\cref{b2:fig:23-cap-tube}把两个小相位条件画成空间宽度与长度。两幅图分别属于频率空间和位置空间；帽的法向高度与管的长度不是同一空间中的两个长度。

\begin{figure}[htbp]
\centering
\begin{tikzpicture}[x=1cm,y=1cm]
\begin{scope}[shift={(-2.8,0)}]
\draw[->] (-1.5,0)--(1.5,0) node[right]{\(\eta\)};
\draw[->] (0,-.3)--(0,1.8) node[above]{法向};
\draw[AnalysisBlue,thick,domain=-1.3:1.3,samples=35] plot(\x,{1.3-.35*\x*\x});
\draw[AnalysisBlue,very thick,domain=-.5:.5,samples=15] plot(\x,{1.3-.35*\x*\x});
\draw[<->] (-.5,.7)--(.5,.7) node[midway,below]{\(\delta\)};
\draw[<->] (.7,1.2125)--(.7,1.3) node[right]{\(\delta^2\)};
\node at (0,-.6){频率帽：切向与法向高度};
\end{scope}
\draw[->] (-.7,.8)--(.6,.8);
\begin{scope}[shift={(2.8,0)}]
\fill[AnalysisBlue!15] (-.4,0) rectangle(.4,1.8);
\draw[AnalysisBlue,thick] (-.4,0) rectangle(.4,1.8);
\draw[<->] (-.4,-.2)--(.4,-.2) node[midway,below]{\(\delta^{-1}\)};
\draw[<->] (.65,0)--(.65,1.8) node[midway,right]{\(\delta^{-2}\)};
\node at (0,2.05){\(x_d\) 法向};
\node at (0,-.85){物理空间相干管};
\end{scope}
\end{tikzpicture}
\caption[频率帽与相干管]{Knapp 尺度的对应，不按比例。移去北极的常相位后，要让 \(x'\cdot\eta\) 与 \(x_dO(|\eta|^2)\) 都小，分别需要 \(|x'|\lesssim\delta^{-1}\) 与 \(|x_d|\lesssim\delta^{-2}\)。细长管的长轴沿帽中心的法向；帽的面积和管的体积仍由正文计算。}
\label{b2:fig:23-cap-tube}
\end{figure}

\subsection{球面尾部的可积性条件}
\label{b2:subsec:230203}

小帽检验只利用一个大而有限的管。要检查无穷远处的积分，取一个不会随尺度收缩的输入更合适。球面上的常数函数便提供了这样的检验。

\Needspace{6\baselineskip}
\begin{proposition}[常密度的严格尾部条件]\label{b2:prop:23-sphere-tail}
设 \(d\geq2\)，\(E\) 是单位球面面积测度的延拓算子。
对于 \(1\leq q<\infty\)，有
\[
 E1\in L^q(\mathbb R^d)
 \quad\Longleftrightarrow\quad
 q>\frac{2d}{d-1}.
\]
因而任何 \(L^s(S^{d-1})\to L^q(\mathbb R^d)\) 的有界性，
若 \(q<\infty\)，都必须满足这个严格不等式。
\end{proposition}
\begin{proof}
球面的旋转不变性使 \(E1(r\omega)\) 只依赖 \(r\)。
相位 \(\omega\cdot\xi\) 在球面上的驻点恰为 \(\omega\) 与 \(-\omega\)，
两处的切向 Hessian 都非退化。按\cref{b2:thm:21-stationary-phase}
分别取两个小驻点邻域，其余部分用非驻相分部积分，可得
\[
 E1(r\omega)=r^{-(d-1)/2}
       \bigl(c_+e^{2\pi\mathrm i r}+c_-e^{-2\pi\mathrm i r}\bigr)
       +O(r^{-(d+1)/2}),\quad r\geq1.
\]
系数 \(c_\pm\) 非零；由球面对称性，它们互为共轭，因而模相等。
旋转不变性也使余项常数不依赖 \(\omega\)。
这给出 \(|E1(r\omega)|\leq Cr^{-(d-1)/2}\)。
原积分在有界区域有界，所以当
\(d-1-q(d-1)/2<-1\) 时，极坐标积分收敛。

反过来，令
\(H(r)=c_+e^{2\pi\mathrm i r}+c_-e^{-2\pi\mathrm i r}\)。
这是非恒零的一周期连续函数。因此存在 \(\kappa>0\) 和
\(A\subset[0,1]\)，使 \(|A|>0\)，且 \(|H(t)|\geq2\kappa\)
对 \(t\in A\) 成立。对充分大的整数 \(n\)，余项在
\(n+A\) 上至多为 \(\kappa r^{-(d-1)/2}\)，从而
\[
 |E1(r\omega)|\geq\kappa r^{-(d-1)/2}
       \quad(r\in n+A).
\]
于是
\[
 \|E1\|_q^q\gtrsim
 \sum_{n\geq n_0} n^{\,d-1-q(d-1)/2}.
\]
这个级数只有在指数小于 \(-1\) 时才收敛。等号时出现调和级数，
也发散。常数函数属于球面上的每个 \(L^s\)，故得到最后的必要条件。
\end{proof}

这里不能把“驻相衰减上界”直接当作发散证明；真正排除端点的是每段径向区间内都有固定比例的下界集合。余弦式领先项确有零点，但这些零点无法消去其余径向区域的质量。

\ProseParagraphBreak
两种必要条件的作用不同：常密度给出严格的水平门槛
\(q>2d/(d-1)\)，小帽给出可含等号的关系
\((d+1)/q\leq(d-1)/s'\)。
在 \(s=2\) 时，小帽条件更强；若允许更大的 \(s\)，两个边界便可能分别起作用。
三维中它们是 \(q>3\) 和 \(4/q\leq2/s'\)。

\subsection{几何障碍}
\label{b2:subsec:230204}

曲率退化时，圆形帽未必反映最强的集中机制。考虑三维中的局部曲面
\[
 S=\{(u,v,u^2+v^4):|u|,|v|<a\},
\]
其中 \(a>0\) 固定，取原生面积测度。对小 \(\delta>0\)，取帽
\[
 |u|\leq\delta^{1/2},\quad |v|\leq\delta^{1/4}.
\]
其高度至多为 \(2\delta\)，面积约为 \(\delta^{3/4}\)。
令 \(g_\delta\) 为该帽的特征函数。若
\[
 |x_1|\leq c\delta^{-1/2},\quad
 |x_2|\leq c\delta^{-1/4},\quad
 |x_3|\leq c\delta^{-1},
\]
则相位 \(2\pi(x_1u+x_2v+x_3(u^2+v^4))\) 的绝对值至多为
\(8\pi c\)。取 \(c\) 足够小，仍可用实部正下界得到
\(|Eg_\delta|\gtrsim\delta^{3/4}\)。
这个长方管的体积约为 \(\delta^{-7/4}\)，所以有界性必须要求
\[
 \delta^{3/4-7/(4q)}
       \lesssim\delta^{3/(4s)},\quad
 \frac7q\leq\frac3{s'}.
\]
特别地，平方输入必须有 \(q\geq14/3\)，比非零曲率曲面的三维指数
\(4\) 更严格。这里曲率在 \(v=0\) 处消失，第四次方方向让帽比二次帽更宽，才产生了更大的相干体积。

\ProseParagraphBreak
这些检验也说明了必要条件的限度。它们都选取了一类特别的输入；通过这些检验，只能说明该类输入没有排除某个指数。任意输入还可能包含许多尺度、位置与方向不同的帽，叠加后的管束相互作用没有包含在单帽计算中。下一节建立对偶与 \(TT^*\) 表述，随后用解析插值证明平方输入的临界估计；更广的目标放在选读部分讨论。
